% Preambulo compartido: paquetes, paleta, estilos de TikZ y macros.
% Responsabilidad unica: configurar el documento. No contiene contenido.

% Tipografia (requiere xelatex; ver README). Pagella para el cuerpo, Heros para
% titulos y diagramas, JetBrains Mono para codigo. Las tres estan instaladas:
% las dos primeras vienen con texlive-fontsrecommended, la tercera es del sistema.
\usepackage{fontspec}
\usepackage{unicode-math}
% Las TeX Gyre viven en texmf y NO estan registradas en fontconfig: se cargan por
% nombre de archivo, no por nombre de familia.
\setmainfont{texgyrepagella}[
  Extension=.otf, UprightFont=*-regular, ItalicFont=*-italic,
  BoldFont=*-bold, BoldItalicFont=*-bolditalic]
\setsansfont{texgyreheros}[
  Extension=.otf, UprightFont=*-regular, ItalicFont=*-italic,
  BoldFont=*-bold, BoldItalicFont=*-bolditalic, Scale=MatchLowercase]
\setmonofont{JetBrains Mono}[Scale=MatchLowercase]
\setmathfont{texgyrepagella-math.otf}
\usepackage[spanish,es-noquoting,es-tabla]{babel}
% expansion no existe en xetex; protrusion si, y es la que mas se nota.
\usepackage[protrusion=true,expansion=false]{microtype}
\usepackage{geometry}
\geometry{a4paper,margin=2.4cm,headsep=12pt}
\usepackage{amsmath}
\usepackage{booktabs}
\usepackage{longtable}
\usepackage{array}
\usepackage{enumitem}
\usepackage{tcolorbox}
\usepackage{tikz}
\usepackage{fancyhdr}
\usepackage{titlesec}
\usepackage[hidelinks]{hyperref}

\usetikzlibrary{positioning,shapes.geometric,arrows.meta,calc,fit,backgrounds}
\tcbuselibrary{skins,breakable}

% --- Paleta -----------------------------------------------------------------
\definecolor{solar}{HTML}{B26B00}   % naranja apagado: fuentes fisicas / energia
\definecolor{tierra}{HTML}{4F6B3A}  % verde: agro / AgroDash
\definecolor{agua}{HTML}{20566E}    % azul: almacenamiento
\definecolor{tinta}{HTML}{1F2328}   % texto
\definecolor{humo}{HTML}{6B7280}    % texto secundario
\definecolor{papel}{HTML}{F5F3EE}   % fondo de cajas
\definecolor{borde}{HTML}{D6D2C8}

\hypersetup{colorlinks=false,pdftitle={AgroVoltaic --- Arquitectura},
            pdfauthor={Proyecto AgroVoltaic (TEC)}}

% --- Titulos ----------------------------------------------------------------
\titleformat{\section}{\Large\bfseries\color{agua}}{\thesection}{0.7em}{}
\titleformat{\subsection}{\large\bfseries\color{tinta}}{\thesubsection}{0.7em}{}
\titlespacing*{\section}{0pt}{2.2ex plus 1ex minus .2ex}{1.2ex plus .2ex}

\pagestyle{fancy}
\fancyhf{}
\fancyhead[L]{\small\color{humo}AgroVoltaic --- Arquitectura del sistema}
\fancyhead[R]{\small\color{humo}\thepage}
\renewcommand{\headrulewidth}{0.4pt}

% --- Macros -----------------------------------------------------------------
\newcommand{\cd}[1]{\texttt{\small #1}}          % identificador / ruta / comando
% \nota: texto secundario. OJO: no admite \\ adentro — el salto de linea de un
% nodo de TikZ es una celda de alineacion y no cruza el grupo. Una \nota por linea.
\newcommand{\nota}[1]{{\color{humo}\small #1}}

\newtcolorbox{clave}[1][]{
  enhanced, breakable, colback=papel, colframe=agua, boxrule=0.6pt,
  left=8pt, right=8pt, top=6pt, bottom=6pt, arc=2pt,
  fonttitle=\bfseries\small, coltitle=white, colbacktitle=agua, #1}

\newtcolorbox{aviso}[1][]{
  enhanced, breakable, colback=white, colframe=solar, boxrule=0.6pt,
  left=8pt, right=8pt, top=6pt, bottom=6pt, arc=2pt,
  fonttitle=\bfseries\small, coltitle=white, colbacktitle=solar, #1}

% --- Estilos de diagrama ----------------------------------------------------
\tikzset{
  comp/.style={
    draw=agua, fill=papel, line width=0.7pt, rounded corners=3pt,
    align=center, inner sep=6pt, minimum height=10mm, font=\small},
  compx/.style={comp, draw=humo, fill=white},          % componente externo
  fuente/.style={comp, draw=solar, fill=solar!8},      % fuente fisica / dato crudo
  agro/.style={comp, draw=tierra, fill=tierra!8},      % region Cartago
  almacen/.style={
    draw=agua, fill=agua!8, line width=0.7pt, cylinder, shape border rotate=90,
    aspect=0.16, align=center, inner sep=6pt, minimum height=13mm, font=\small},
  etapa/.style={comp, minimum width=27mm},
  flecha/.style={-{Stealth[length=2.4mm]}, line width=0.7pt, draw=tinta},
  flechad/.style={flecha, dashed, draw=humo},
  etiq/.style={font=\scriptsize\itshape, color=humo, inner sep=2pt, fill=white},
  grupo/.style={draw=borde, line width=0.6pt, rounded corners=4pt, inner sep=8pt},
  % Sin guiones dentro de los diagramas: partir "ana-lizador" en una caja de
  % 27 mm se lee peor que dejar la caja un poco mas holgada.
  every picture/.append style={execute at begin picture={%
    \hyphenpenalty=10000\relax\exhyphenpenalty=10000\relax}},
  titgrupo/.style={font=\scriptsize\bfseries, color=humo},
}
